% !TeX root = ../../Book3.tex
\section{Picard 小定理}
\label{b3:sec:1404}

Picard 小定理把省略值从族的紧性转成单个整函数的刚性。圆盘覆盖将目标限制变成一个有界整函数，Liouville 遂可直接使用。

\subsection{覆盖空间证明}
\label{b3:subsec:140401}

\term[picard-xiao-dingli]{Picard 小定理}[little Picard theorem]限制整函数的省略值。

\begin{theorem}[Picard 小定理]\label{b3:thm:14-picard-small}
非恒常整函数至多省略一个有限复值；非恒常全平面亚纯函数至多省略两个球面值。
\end{theorem}
\begin{proof}
若亚纯函数省略三个不同球面值，它是从单连通平面到三穿孔球面的全纯映射。路径提升给到圆盘万有覆盖的全纯提升，Liouville 使提升常值，所以原函数常值。整函数本来不取球面值无穷远，再省略两个有限值就省略三个球面值，因此也不可能。
\end{proof}
这项证明使用已声明的深层覆盖模型输入，提升与 Liouville 部分完整。它不将“三穿孔球面有圆盘覆盖”从 Picard 定理反推，避免循环。

\subsection{省略值}
\label{b3:subsec:140402}

\(e^z\) 省略0，说明整函数的一值例外确会发生。作为亚纯函数它同时省略0与无穷远，达到球面版本的两值上界。多项式则不省略任何有限值，代数基本定理已经给出其满射性。

\ProseParagraphBreak
“至多一个”不表示必有一个。\(\sin z\) 对每个 \(a\in\mathbb C\) 都取值：令 \(w=e^{\mathrm{i}z}\)，方程 \((w-w^{-1})/(2\mathrm{i})=a\) 化为常数项为 \(-1\) 的二次方程，根非零，再取复对数求 \(z\)。小定理只控制遗漏，不直接计数原像。

\subsection{应用}
\label{b3:subsec:140403}

若非常值整函数 \(f,g\) 满足 \(e^f+e^g=1\)，则 \(e^f\) 同时不取0与1，Picard 小定理迫使它常值，求导再给 \(f\) 常值，与假设矛盾。因此不存在这样的非常值整函数对。相同方法可检验一些指数型函数方程：先把式子变成两个省略值，再用刚性，而非展开无穷系数逐项比较。

\ProseParagraphBreak
区域换成圆盘后这种结论就不成立。常值附近足够小的全纯扰动可同时避开0与1。整平面的尺度无限是 Liouville 提升证明中的决定条件，不能把小定理搬到任意区域。

\ProseParagraphBreak
对于超越整函数，无穷远是本性奇点。理由是把 \(F(z)\) 改写成 \(F(1/t)\)：若在0可去，则 \(F\) 在大圆外有界，再与圆内紧性合并，由 Liouville 为常数；若为极点，则 \(F\) 有多项式增长，Cauchy 导数估计迫使充分高阶 Taylor 系数为0，故为多项式。超越整函数既非前者也非后者，所以大 Picard 给它在每个大圆外取遍所有复值至多除去一个，而且其余值取无限次。

\ProseParagraphBreak
小 Picard 只限制全局省略值，大 Picard 还控制重复次数。\(e^z\) 省略0，其他每个非零值在任意大圆外都有无穷多个对数原像。\(\sin z\) 不省略任何复值，利用 \(e^{\mathrm{i}z}\) 的二次方程及周期也可显式看到无限原像。这些例子与定理一致，但不能由几个常见超越函数的行为代替一般本性奇点的证明。
